\chapter{Matrix-Valued Tight-Binding Systems}

\section{Introduction}

The one-dimensional chain studied in the previous chapters possesses a
single quantum amplitude at each lattice site. Consequently, the
wavefunction is represented by a scalar quantity,

\[
	\psi_j.
\]

Although this description is sufficient for a simple one-dimensional
chain, most physical systems encountered in condensed matter physics
require several amplitudes to describe the local quantum state.

Examples include

\begin{itemize}

	\item graphene, whose unit cell contains two inequivalent sublattices;

	\item bilayer graphene, where the wavefunction contains four amplitudes;

	\item spin-dependent systems, in which each orbital possesses two spin
	      components;

	\item superconducting systems described by the
	      Bogoliubov--de Gennes formalism, where electron and hole degrees of
	      freedom are simultaneously present.

\end{itemize}

Despite their apparent diversity, all these systems share the same
mathematical structure.

The objective of this chapter is to develop a general tight-binding
formalism capable of describing any nearest-neighbor lattice whose local
wavefunction contains an arbitrary number of internal degrees of
freedom.

The resulting equations constitute the mathematical foundation of the
remaining chapters of this book.

\section*{Notation Summary}

Throughout this chapter, the following notation will be adopted.

\begin{center}
	\renewcommand{\arraystretch}{1.3}
	\begin{tabular}{ll}
		\hline
		Symbol          & Meaning                                                 \\
		\hline
		$j$             & Unit-cell index along the transport direction           \\
		$N$             & Dimension of the local Hilbert space                    \\
		$\Psi_j$        & Wavefunction associated with unit cell $j$              \\
		$\Phi_j$        & Augmented state vector used in the transfer formulation \\
		$H_j$           & Onsite Hamiltonian of unit cell $j$                     \\
		$T$             & Hopping matrix coupling neighboring unit cells          \\
		$E$             & Incident electron energy                                \\
		$I$             & Identity matrix of dimension $N\times N$                \\
		$\mathcal{M}_j$ & Local transfer matrix                                   \\
		$\mathcal{M}$   & Total transfer matrix                                   \\
		\hline
	\end{tabular}
\end{center}

\section{Local Hilbert Space}

Instead of associating a single complex number with each lattice site,
we now associate an $N$-component vector,

\begin{equation}
	\Psi_j
	=
	\begin{pmatrix}
		\psi_{1,j} \\
		\psi_{2,j} \\
		\vdots     \\
		\psi_{N,j}
	\end{pmatrix},
	\label{eq:vector-wavefunction}
\end{equation}
where $N$ denotes the dimension of the local Hilbert space.

Each component may represent a physical degree of freedom such as

\begin{itemize}

	\item sublattice,

	\item orbital,

	\item spin,

	\item layer,

	\item electron-hole sector,

	\item or any combination of these.

\end{itemize}

Consequently, the one-dimensional chain studied previously corresponds
simply to the particular case

\[
	N=1.
\]

Rather than representing a different physical problem, the scalar
formalism of Chapters~1 and~2 now appears as the simplest member of a
much more general family of tight-binding models.

%%%%%%%%%%%%%%%%%%%%%%%%%%%%%%%%%%%%%%%%%%%%%%%%%%%%%%%%%%

\begin{center}
	\fbox{
		\begin{minipage}{0.90\linewidth}

			\textbf{Physical Insight}

			\vspace{1ex}

			The lattice index $j$ specifies \emph{where} the particle is located.

			The vector index specifies \emph{which internal quantum state} the
			particle occupies.

			The transfer matrix propagates the wavefunction along the lattice,
			whereas the internal matrices describe the local quantum dynamics.

		\end{minipage}
	}
\end{center}

%%%%%%%%%%%%%%%%%%%%%%%%%%%%%%%%%%%%%%%%%%%%%%%%%%%%%%%%%%

\section{General Tight-Binding Hamiltonian}

The scalar onsite energy introduced in the previous chapters is replaced
by an $N\times N$ Hermitian matrix,

\[
	H_j,
\]
which describes the local Hamiltonian associated with lattice site $j$.

Likewise, the scalar hopping parameter becomes an arbitrary coupling
matrix,

\[
	T,
\]
connecting neighboring lattice sites.

The Schrödinger equation therefore assumes the general form

\begin{equation}
	(EI-H_j)\Psi_j
	+
	T^\dagger\Psi_{j-1}
	+
	T\Psi_{j+1}
	=
	0.
	\label{eq:general-tight-binding}
\end{equation}

Equation~(\ref{eq:general-tight-binding}) is the central equation of the
present chapter.

Every tight-binding model considered throughout this book will be
obtained from Eq.~(\ref{eq:general-tight-binding}) by specifying the
matrices $H_j$ and $T$.

%%%%%%%%%%%%%%%%%%%%%%%%%%%%%%%%%%%%%%%%%%%%%%%%%%%%%%%%%%

\begin{center}
	\fbox{
		\begin{minipage}{0.92\linewidth}

			\textbf{Key Equation}

			\vspace{1ex}

			The matrix equation

			\[
				(EI-H_j)\Psi_j
				+
				T^\dagger\Psi_{j-1}
				+
				T\Psi_{j+1}
				=
				0
			\]

			constitutes the most general nearest-neighbor tight-binding equation
			considered in this book.

			The remaining chapters differ only in the explicit forms of the onsite
			Hamiltonian $H_j$ and the coupling matrix $T$.

		\end{minipage}
	}
\end{center}



\section{General Matrix Tight-Binding Equation}

Equation~(\ref{eq:general-tight-binding}) introduces two matrix-valued
objects that completely characterize the quantum system:

\begin{itemize}

	\item the onsite Hamiltonian,
	      $H_j$;

	\item the hopping matrix,
	      $T$.

\end{itemize}

Their physical roles are fundamentally different.

The onsite Hamiltonian describes the quantum dynamics inside a single
lattice cell, whereas the hopping matrix describes the coupling between
neighboring cells.

This distinction is illustrated schematically as

\[
	\boxed{
		\Psi_{j-1}
		\quad
		\overset{T^\dagger}{\longleftrightarrow}
		\quad
		\Psi_j
		\quad
		\overset{T}{\longleftrightarrow}
		\quad
		\Psi_{j+1}.
	}
\]

The local Hamiltonian acts only inside each unit cell,

\[
	H_j :
	\Psi_j
	\longrightarrow
	\Psi_j,
\]

while the hopping matrices connect neighboring cells,

\[
	T :
	\Psi_j
	\longrightarrow
	\Psi_{j+1},
\]

and

\[
	T^\dagger :
	\Psi_j
	\longrightarrow
	\Psi_{j-1}.
\]

Consequently, Eq.~(\ref{eq:general-tight-binding})

\[
	(EI-H_j)\Psi_j
	+
	T^\dagger\Psi_{j-1}
	+
	T\Psi_{j+1}
	=
	0
\]

should be interpreted as a balance between three distinct physical
processes:

\begin{enumerate}

	\item
	      the local quantum evolution inside the unit cell;

	\item
	      the coupling with the previous unit cell;

	\item
	      the coupling with the next unit cell.

\end{enumerate}

Although these processes acquire different physical interpretations in
graphene, bilayer graphene and superconducting systems, their
mathematical structure remains exactly the same.

%%%%%%%%%%%%%%%%%%%%%%%%%%%%%%%%%%%%%%%%%%%%%%%%%%%%%%%%%%

\begin{center}
	\fbox{
		\begin{minipage}{0.92\linewidth}

			\textbf{Physical Insight}

			\vspace{1ex}

			The onsite Hamiltonian determines the internal quantum dynamics of a
			unit cell.

			The hopping matrix determines how quantum information propagates from
			one unit cell to the next.

			In the transfer-matrix formulation, propagation occurs along the lattice
			index, whereas the matrices describe the internal quantum structure of
			each unit cell.

		\end{minipage}
	}
\end{center}

\section{Transfer-Matrix Formulation}

\subsection{Recurrence Relation}
The objective of this section is to rewrite the matrix-valued
tight-binding equation as a first-order recurrence relation, exactly as
was done for the scalar one-dimensional chain in Chapter~2.

The starting point is the general nearest-neighbor tight-binding
equation,

\begin{equation}
	(EI-H_j)\Psi_j
	+
	T^\dagger\Psi_{j-1}
	+
	T\Psi_{j+1}
	=
	0.
	\label{eq:general-tb}
\end{equation}

This equation couples three neighboring unit cells. In order to obtain a
transfer-matrix formulation, the wavefunction in the next cell must be
written explicitly in terms of the two previous ones.

Rearranging Eq.~\eqref{eq:general-tb}, we obtain

\begin{equation}
	T\Psi_{j+1}
	=
	-
	(EI-H_j)\Psi_j
	-
	T^\dagger\Psi_{j-1}.
	\label{eq:isolate-forward}
\end{equation}

Assuming that the hopping matrix $T$ is invertible, multiplication by
$T^{-1}$ gives

\begin{equation}
	\Psi_{j+1}
	=
	-
	T^{-1}(EI-H_j)\Psi_j
	-
	T^{-1}T^\dagger\Psi_{j-1}.
	\label{eq:matrix-recurrence}
\end{equation}

Equation~\eqref{eq:matrix-recurrence} is the direct generalization of
the recurrence relation obtained for the scalar one-dimensional chain.

The difference is entirely contained in the replacement of scalar
coefficients by matrices acting on the internal Hilbert space.

%%%%%%%%%%%%%%%%%%%%%%%%%%%%%%%%%%%%%%%%%%%%%%%%%%%%%%%%%%

\begin{center}
	\fbox{
		\begin{minipage}{0.90\linewidth}

			\textbf{Physical Insight}

			\vspace{1ex}

			Equation~\eqref{eq:matrix-recurrence} shows that the wavefunction in unit
			cell $j+1$ depends only on the wavefunctions in the two preceding unit
			cells.

			This is precisely the same propagation principle encountered in the
			scalar one-dimensional chain. The only difference is that each
			wavefunction now possesses an internal vector structure.

		\end{minipage}
	}
\end{center}


%%%%%%%%%%%%%%%%%%%%%%%%%%%%%%%%%%%%%%%%%%%%%%%%%%%%%%%%%%
\subsection{Augmented State Vector}

Equation~\eqref{eq:matrix-recurrence} is a second-order recurrence
relation because the wavefunction in unit cell $j+1$ depends on the
wavefunctions in the two preceding unit cells.

As in Chapter~2, the recurrence can be converted into a first-order
matrix equation by enlarging the state vector.

We therefore define the augmented state vector

\begin{equation}
	\Phi_j
	=
	\begin{pmatrix}
		\Psi_j \\
		\Psi_{j-1}
	\end{pmatrix},
	\label{eq:augmented-state}
\end{equation}

which contains the wavefunctions associated with two neighboring unit
cells.

The corresponding vector in the next propagation step is

\begin{equation}
	\Phi_{j+1}
	=
	\begin{pmatrix}
		\Psi_{j+1} \\
		\Psi_j
	\end{pmatrix}.
	\label{eq:augmented-next}
\end{equation}

Substituting Eq.~\eqref{eq:matrix-recurrence} into the first component of
Eq.~\eqref{eq:augmented-next} yields

\begin{equation}
	\Phi_{j+1}
	=
	\begin{pmatrix}
		-T^{-1}(EI-H_j)\Psi_j
		-
		T^{-1}T^\dagger\Psi_{j-1}
		\\[1ex]
		\Psi_j
	\end{pmatrix}.
	\label{eq:augmented-intermediate}
\end{equation}

The first component contains the recurrence relation derived in the
previous subsection, while the second component simply stores the current
wavefunction for use during the next propagation step.

Equation~\eqref{eq:augmented-intermediate} is now ready to be written in
matrix form.



\subsection{Augmented State Vector}

Equation~\eqref{eq:matrix-recurrence} is a second-order recurrence
relation because the wavefunction in unit cell $j+1$ depends on the
wavefunctions in unit cells $j$ and $j-1$.

As in the scalar problem discussed in Chapter~2, this recurrence relation
can be written as a first-order matrix equation by enlarging the state
vector.

We define the augmented state vector as

\begin{equation}
	\Phi_j
	======

	\begin{pmatrix}
		\Psi_j \
		\Psi_{j-1}
	\end{pmatrix}.
	\label{eq:augmented-state-vector}
\end{equation}

Since each vector $\Psi_j$ has $N$ components, the augmented vector
$\Phi_j$ has $2N$ components.

At the next propagation step,

\begin{equation}
	\Phi_{j+1}
	\begin{pmatrix}
		\Psi_{j+1} \
		\Psi_j
	\end{pmatrix}.
	\label{eq:next-augmented-state-vector}
\end{equation}

Using the recurrence relation

\begin{equation}
	\Psi_{j+1}

	T^{-1}(EI-H_j)\Psi_j

	T^{-1}T^\dagger\Psi_{j-1},
	\label{eq:matrix-recurrence-repeated}
\end{equation}

we obtain

\begin{equation}
	\Phi_{j+1}
	\begin{pmatrix}

		## T^{-1}(EI-H_j)\Psi_j

		T^{-1}T^\dagger\Psi_{j-1}
		[1ex]
		\Psi_j
	\end{pmatrix}.
	\label{eq:augmented-intermediate}
\end{equation}

The upper component computes the wavefunction in the next unit cell.
The lower component stores the current wavefunction so that it becomes
the previous-cell wavefunction during the next propagation step.

%%%%%%%%%%%%%%%%%%%%%%%%%%%%%%%%%%%%%%%%%%%%%%%%%%%%%%%%%%

\subsection{General Local Transfer Matrix}

Equation~\eqref{eq:augmented-intermediate} can be written as a block
matrix equation:

\begin{equation}
	\begin{pmatrix}
		\Psi_{j+1} \
		\Psi_j
	\end{pmatrix}
	\begin{pmatrix}

		T^{-1}(EI-H_j)
		 &
		-

		T^{-1}T^\dagger
		[1ex]
		I
		 &
		0
	\end{pmatrix}
	\begin{pmatrix}
		\Psi_j \
		\Psi_{j-1}
	\end{pmatrix}.
	\label{eq:block-transfer-equation}
\end{equation}

Using the definition of the augmented state vector, the propagation
equation becomes

\begin{equation}
	\Phi_{j+1}
	\mathcal{M}_j\Phi_j,
	\label{eq:general-transfer-relation}
\end{equation}

where the generalized local transfer matrix is

\begin{equation}
	\mathcal{M}_j
	\begin{pmatrix}

		T^{-1}(EI-H_j)
		 &
		-

		T^{-1}T^\dagger
		[1ex]
		I
		 &
		0
	\end{pmatrix}.
	\label{eq:general-local-transfer-matrix}
\end{equation}

The matrix $\mathcal{M}_j$ has dimension $2N \times 2N$. Each block has
dimension $N \times N$.

The upper-left block,

\begin{equation}
	A_j
	T^{-1}(EI-H_j),
	\label{eq:transfer-block-a}
\end{equation}

contains the energy dependence and the onsite Hamiltonian of unit cell
$j$.

The upper-right block,

\begin{equation}
	B
	=
	T^{-1}T^\dagger,
	\label{eq:transfer-block-b}
\end{equation}

contains the contribution from the previous unit cell.

The lower blocks,

\begin{equation}
	\begin{pmatrix}
		I & 0
	\end{pmatrix},
\end{equation}

store the current wavefunction in the second part of the augmented state
vector.

%%%%%%%%%%%%%%%%%%%%%%%%%%%%%%%%%%%%%%%%%%%%%%%%%%%%%%%%%%

\begin{center}
	\fbox{
		\begin{minipage}{0.90\linewidth}

			\textbf{Framework Equation}

			\vspace{1ex}

			The generalized local transfer matrix is

				[
					\mathcal{M}_j

					\begin{pmatrix}

						T^{-1}(EI-H_j)
						 &
						-

						T^{-1}T^\dagger
						[1ex]
						I
						 &
						0
					\end{pmatrix}.
				]

			This matrix propagates the augmented state vector according to

				[
					\Phi_{j+1}
					==========

					\mathcal{M}_j\Phi_j.
				]

			All matrix-valued nearest-neighbor systems considered later will be
			obtained by specifying the onsite Hamiltonian $H_j$ and the hopping
			matrix $T$.

		\end{minipage}
	}
\end{center}

%%%%%%%%%%%%%%%%%%%%%%%%%%%%%%%%%%%%%%%%%%%%%%%%%%%%%%%%%%

\subsection{Recovery of the Scalar Limit}

The scalar one-dimensional chain must be recovered when the local Hilbert
space has dimension

\begin{equation}
	N=1.
\end{equation}

In this case,

\begin{equation}
	H_j=\epsilon_j,
	\qquad
	T=t,
	\qquad
	T^\dagger=t,
	\qquad
	I=1.
\end{equation}

Substitution into Eq.~\eqref{eq:general-local-transfer-matrix} gives

\begin{equation}
	\mathcal{M}_j
	\begin{pmatrix}
		-\dfrac{E-\epsilon_j}{t}
		 &
		-1
		[1ex]
		1
		 &
		0
	\end{pmatrix}.
	\label{eq:scalar-transfer-limit}
\end{equation}

This is precisely the local transfer matrix derived for the scalar chain
in Chapter~2.

The scalar formulation is therefore recovered as the special case
$N=1$ of the matrix-valued theory.

%%%%%%%%%%%%%%%%%%%%%%%%%%%%%%%%%%%%%%%%%%%%%%%%%%%%%%%%%%

\begin{center}
	\fbox{
		\begin{minipage}{0.90\linewidth}

			\textbf{Validation Principle}

			\vspace{1ex}

			The matrix-valued implementation must reproduce the previously validated
			scalar solver when the internal dimension is set to $N=1$.

			This limiting case will provide the first automated test of the
			generalized transfer-matrix solver.

		\end{minipage}
	}
\end{center}
